\section{Cross-Layer Integrity in Hybrid Web Environments}
\label{sec:setting}
\subsection{Multiple Views of One Execution}
An Android hybrid Web execution provides three complementary, partially overlapping views. Native supplies system and hardware observations; the WebView Host supplies container, configuration, provider, and layout context; App Web exposes JavaScript-visible reports inside that WebView~\cite{android2026webview}. Let $E$ denote the relevant runtime state and container context, beyond a static device model or identifier. We write these observations as $N(E)$, $H(E)$, and $W(E)$ without assuming noiseless, complete, or deterministic API outputs. Associating observations from one execution supports comparison, without implying atomic sampling, trusted collection, or hardware attestation.

The implemented \appfp measurement contract catalogs 84 Native, 26 Host, and 67 App Web fields. An external browser runs a separate Web execution on the same device. \browserfp catalogs 67 fields from that execution; \pairedfp represents linked App--browser observations. These counts describe implementation catalogs, not independent information or the inputs used by every detector.

\subsection{Cross-Layer Consistency as Integrity Evidence}
Integrity here concerns report manipulation; available evidence limits which relations can be assessed. A local Web-report modification can be illustrated as
\[
\begin{aligned}
x_E &= (N(E),H(E),W(E)),\\
x'_E &= (N(E),H(E),W'(E)).
\end{aligned}
\]
Each semantic relation $C_j$ constrains a relevant subset of observations, such as Native--Web or Host--Web. Its interpretation requires an applicability domain, usable operands, and the necessary measurement context. A difference provides evidence only insofar as it violates that measurement constraint.

Local editing may break an applicable relation when other observations still provide distinguishing context: a plausible Web value can then be incompatible with its reference. Some edits remain within compatible ranges. A normal system timezone change may instead move observations together while preserving their relation. Sampling times and implementation differences still limit this interpretation. More views alone do not guarantee a gain; useful evidence depends on informative references and interpretable constraints. This is a mechanism examined empirically, not a necessary-and-sufficient characterization.

\hg evaluates concrete checks using their current required observations. $E$ is explanatory, not an available ground-truth label: the detector neither reconstructs $E$ nor proves that all observations describe a single authentic environment. The notation does not imply access to pre-change baselines, recovery records, or experiment-stage labels. $C_j$ describes compatibility, whereas $r_j(x)$ in Section~\ref{sec:method} denotes a detection condition's state.

The three-view model does not require every relation to use all views or imply a triangle of validated checks. \full combines within-view and cross-view checks; Host geometry remains a separate, unintegrated study; no standalone Native--Host detection rule has been validated.

\subsection{Adversary and Scope of Integrity Claims}
The tested manipulations alter selected resource, browser, locale, display, and graphics reports through controlled tools, browser emulation, or runtime overrides, without necessarily changing the physical device. Only a measured effect at the target surface qualifies as an \emph{effective modification}; operations with no observable effect remain separate.

Native and Host can also be manipulated. A reference left unchanged by a tested intervention is neither a trusted root nor evidence of statistical or administrative independence. Coordinated changes across compared views, collector control, or backend compromise may evade checks; resistance to these capabilities is unproven. Such evasion differs from misses caused by limited implemented relations or observation coverage.

An alarm means an applicable integrity-check condition triggered; it does not confirm fraud, account takeover, malicious intent, or attacker identity. No alert does not establish authenticity. Evaluation units are saved \emph{record positions} or applicable paired positions; repeated operations are not independent devices. Experiment roles and operation records serve evaluation, not prediction.

\subsection{Normal Variation and Missing Evidence}
Normal system timezone, language, orientation, and application-layout changes affect comparability, as do API precision, units, support, and container state. Semantic relations can account for unit conversion, coarse reporting, measurement intervals, and container applicability. Legitimate differences require interpretation within these constraints; they are not automatically violations or unknown outcomes.

Observation availability, a relation among observations, and the detector's final output are distinct. For a selected condition, \hgstate{T} marks a triggered alarm condition; \hgstate{F} means it does not trigger, without certifying authenticity; \hgstate{U} means required evidence is unavailable or insufficient, or the relation is inapplicable. Execution or association failure is recorded separately as \hgstate{FAILED}. An empty model also has a separate status, not an ordinary normal result.

A missing required operand can prevent judging its relation. An unused field's absence does not independently fail the detector, and a check using only some views does not require all three to be complete. Missing evidence is not filled in as normal. Section~\ref{sec:method} explains how \hg implements these conditional relations and evidence states as executable checks and detectors.
